\chapter{Fejlesztői dokumentáció}
\label{ch:impl}

Ez a fejezet az alkalmazás elemzését, tervezését, implementációját és tesztelését mutatja be.

%% ============================================================
\section{Elemzés}
%% ============================================================

\subsection{Feladatleírás}

A feladat egy olyan interaktív alkalmazás fejlesztése, amely:
\begin{compactenum}
\item Modellezi a Rush Hour tologatós rejtvényt~\cite{rushhour_puzzle} egy immutábilis állapot-átmenet rendszerként egy $6\times6$-os táblán, $1\times2$-es autókkal és $1\times3$-as buszokkal.
\item Implementál három konfigurálható keresőalgoritmust (tabu keresés, visszalépéses keresés, A$^C$ gráfkeresés) négy heurisztikus függvénnyel.
\item Vizualizálja a feltárt állapotgráfot 3D-ben, erőalapú elrendezést (force-directed layout) és Barnes–Hut optimalizációt~\cite{barneshut1986} használva.
\item Lehetővé teszi a kézi játékot és az állapotok interaktív vizsgálatát a gráfcsúcsokra való kattintással.
\end{compactenum}

\subsection{Funkcionális leírás}

A funkcionális követelményeket a \ref{tab:userstory}. táblázatban található felhasználói történetek (user stories) foglalják össze.

\begin{center}
\begin{longtable}{ | p{0.20\textwidth} | p{0.20\textwidth} | p{0.20\textwidth} | p{0.30\textwidth} | }

\hline
\multicolumn{4}{|c|}{\textbf{Felhasználói történetek}}
\\ \hline
\textbf{Cél} & \textbf{Előfeltétel} & \textbf{Esemény} & \textbf{Eredmény}
\\ \hline \hline
\endfirsthead

\hline
\textbf{Cél} & \textbf{Előfeltétel} & \textbf{Esemény} & \textbf{Eredmény}
\\ \hline \hline
\endhead

\hline
\endfoot
\endlastfoot
Kiválasztani egy pályát & A feladvány választó fül meg van nyitva & Rányomok a megfelelő pálya előnézet rublikájára & Betöltődik a pálya kezőállapota, a állapotgráfon csak a kezdőállapot látszik, átvált a jobb panel az algoritmus beállítások fülre \\ \hline
Felfedezni az állapotgráf csúcsait & Kereső algoritmus beállítások fül meg van nyitva & Rányomok a csúcsfelfedezés gombjára & Az állapotgráfon felvesszük a felfedezett csúcsokat, a célcsúcsok ki vannak emelve \\ \hline
Megállítani a csúcsfelfedezést & Épp fut csúcsfelfedezés & Rányomok a csúcsfelfedezés megállítása gombjára & A csúcsfelfedezés megáll \\ \hline
Kiválasztani milyen keresőalgoritmust legyen kiválasztva & Kereső algoritmus beállítások fül meg van nyitva & Kiválasztom a megfelelő keresőalgoritmust az opciók közül & A keresőalgoritmus beállításra kerül \\ \hline
Kiválasztani milyen heurisztikus függvény legyen kiválasztva & Kereső algoritmus beállítások fül meg van nyitva & Kiválasztom a megfelelő heurisztikus függvényt az opciók közül & A heurisztika beállításra kerül \\ \hline
Futtatni a beállított keresőalgoritmust & Kereső algoritmus beállítások fül meg van nyitva & Rányomok a keresőalgoritmus alkalmazása gombra & A kiválasztott beállítások alapján alkalmazva lesz a kereső algoritmus a betöltött feladványra \\ \hline
Elindítani/ folytatni a keresőalgoritmust & A játéktábla fül meg van nyitva & Rányomok a keresőalgoritmus indítása/folytatása gombjára & A keresőalgoritmus elindul, a gráfon ki vannak emelve a globális munkaterületen tárolt csúcsok és élek, ha a keresőalgoritmus talál célállapotot, akkor leáll az algoritmus \\ \hline
Megállítani a keresőalgoritmust & A játéktábla fül meg van nyitva és épp fut keresőalgoritmus & Rányomok a keresőalgoritmus megállítása gombjára & A keresőalgoritmus megáll, a gráfon a globális munkaterületen tárolt csúcsok és élek ki vannak emelve \\ \hline
Lépésenként futtatni a keresőalgoritmust & A játéktábla fül meg van nyitva & Rányomok a keresőalgoritmus lépésenkénti futtatása gombjára & A keresőalgoritmus egy lépést futtat, a gráfon ki vannak emelve a globális munkaterületen tárolt csúcsok és élek \\ \hline
Beállítani a keresőalgoritmus sebességét & A játéktábla fül meg van nyitva & A sebesség állító csúszkát a megfelelő helyre húzom & A keresőalgoritmus minden lépése a beállított sebességgel történik \\ \hline
Manuálisan lépni a táblán & A játéktábla fül meg van nyitva & A játéktáblán kattintok a lépésnek megfelelő piros nyílra & A játéktáblán elmozdul az autó, az állapotgráfon felfedezzük, kitüntetjük az új állapot csúcsát \\ \hline
A játéktáblát manuális módba váltani & A játéktábla fül meg van nyitva & A játéktábla fölötti manuális gombra kattintok & A játéktábla manuális módba vált, az állapotgráfon ki van tüntetve a megjelenített állapot, ha futott algoritmus az a háttérben fut tovább \\ \hline
Az algoritmus lépéseit megjeleníteni a táblán & A játéktábla fül meg van nyitva & A játéktábla fölötti algoritmus gombra kattintok & A játéktábla algoritmus módba vált, a táblán az algoritmus által legutóbb kiterjesztett állapot van megjelenítve, a játéktábla mutatja az algoritmus újjab kiterjesztett állapotait \\ \hline
Egy felfedezett állapotba ugrani & Egy feladvány be van töltve & Rányomok (M1) egy felfedezett állapot csúcsára az állapotgráfon & A játéktábla a kiválasztott állapotot mutatja, a gráfon a kiválasztott csúcs ki van emelve \\ \hline
Forgarni az állapotgráfot & Egy feladvány be van töltve & Az állapotgráfon kattintok (M2) és húzom az egeret & Az állapotgráf forog az aktív csúcs körül az egér mozgásának megfelelően \\ \hline
Mozgatni az állapotgráfot & Egy feladvány be van töltve & A WASD gombok közül nyomom valamelyiket & Az állapotgráf elmozdul a vásznon a gombnak megfelelően \\ \hline
Az állapotgráfon elrejteni/megjeleníteni a nem releváns éleket, csúcsokat & A játéktábla fül meg van nyitva & Rákattintok az irreleváns élek és csúcsok elrejtésének kapcsolójára & Az állapotgráfon elrejtődnek/megjelennek az irreleváns élek, csúcsok \\ \hline
Követni az állapotgráfon az algoritmus által kiterjesztett csúcsokat & A játéktábla fül meg van nyitva & Rákattintok az algoritmus csúcsai követésének kapcsolójára & Az állapotgráf kamerája követi az algoritmus által kiterjesztett csúcsokat \\ \hline
Nagyítani/ kicsinyíteni az állapotgráfot & Egy feladvány be van töltve & Az állapotgráfon az egér görgőjét használom & Az állapotgráf nagyítódik/kicsinyítódik az egér görgőjének megfelelően \\ \hline
Kilépni a programból & A program fut & Bezárom a program ablakát & A program bezárul \\ \hline
\caption{Felhasználói történetek}
\label{tab:userstory}
\end{longtable}
\end{center}


\subsection{Nem-funkcionális követelmények}

\begin{compactitem}
\item \textbf{Teljesítmény}: a 3D gráfnak interaktívnak kell maradnia akár több ezer csúcs esetén is. A Barnes–Hut optimalizációnak kezelhető szinten kell tartania a képkockánkénti számítási igényt.
\item \textbf{Hordozhatóság}: önálló futtatható fájlok Windowsra és Linuxra; fordíthatóság más platformokra a Godot szerkesztőn keresztül.
\item \textbf{Használhatóság}: intuitív felhasználói felület színkódolt gráfelemekkel és reagáló vezérlőkkel.
\item \textbf{Helyesség}: a keresőalgoritmusoknak az elvárt tulajdonságok szerint kell működniük (pl. az A$^C$ gráfkeresésnek monoton heurisztikával optimális megoldást kell találnia).
\end{compactitem}

%% ============================================================
\section{Tervezés}
%% ============================================================

\subsection{Szoftverarchitektúra}

% TODO cite something?
Az alkalmazás a \textbf{Modell--Nézet (Model--View)} architektúrát követi:

\begin{description}
\item[Modell réteg] (\texttt{src/Model/}) Tiszta C\# osztályokat tartalmaz a játékállapot reprezentációjára, a lépésgenerálásra, a heurisztikus kiértékelésre és a keresőalgoritmusokra. A modell nem függ a Godot motortól, egyedül a Godot-ban definiált \texttt{Vector2I} struktúrát használja, hogy ne kelljen később konvertálni azt.
\item[Nézet réteg] (\texttt{src/Nodes/}) Olyan Godot node scripteket tartalmaz, amelyek a motor alaposztályaiból származnak (\texttt{Control}, \texttt{RigidBody3D}, \texttt{MeshInstance3D}, stb.). Ezek az osztályok kezelik a kirajzolást, a fizikai szimulációt, a felhasználói felületet és bevitelt. Feliratkoznak a modell eseményeire (\texttt{NewEdge}, \texttt{PathChange}, \texttt{NewCurrent}, \texttt{Terminated}), hogy reagáljanak az algoritmus előrehaladására.
\end{description}

A modell és a nézet közötti kommunikáció az absztrakt \texttt{RHSolver} alaposztályon definiált C\# eseményeken (events) keresztül valósul meg. A nézet feliratkozik ezekre az eseményekre, és létrehozza vagy frissíti a megfelelő vizuális elemeket (csúcsok, élek, kiemelések).

\subsection{Modell}
\label{sec:model}

A modell fontosabb részeinek az osztálydiagramja a \ref{fig:model-state-class}. és \ref{fig:model-solvers-class}. ábrán látható.

\begin{figure}[H]
\centering
\includegraphics[width=1.0\textwidth]{model_state}
\caption{Az állapot ábrázolás osztálydiagramja}
\label{fig:model-state-class}
\end{figure}

Az állapotok és bábuk mind módosíthatatlan osztályok.


\begin{description}
\item[\texttt{RHGameState}] Egy táblakonfiguráció módosíthatatlan reprezentációja, a módosíthatatlanságot az \texttt{init;} kulcsszóval és a \textit{readonly} adat struktúrákkal biztosítjuk. Az elhelyezett járműveket \texttt{PlacedRHPiece} objektumok tömbjét tárolja, ezekből pedig generál egy $6\times6$-os egész értékű mátrixot tárol a gyors mező eléréshez. Biztosítja a \texttt{GetPossibleMoves()}, \texttt{WithMove()}, \texttt{IsSolved()} metódusokat, valamint érték-alapú egyenlőségvizsgálatot amit hash kódokkal van felgyorsítva. Mivel nem módosítható az adatszerkezet, ezért elég csak inicializáláskor egyszer kiszámolni a hash kódját az állásra, és később mindig ezt adja vissza.

\item[\texttt{PlacedRHPiece}] Összekapcsol egy \texttt{RushHourPiece}-t (Car, MainCar vagy Bus) egy táblapozícióval és iránnyal. A \texttt{WithMove()} metódus egy új példányt ad vissza frissített pozícióval.

\item[\texttt{Move} / \texttt{StateMove}] A \texttt{Move} egy önmagában nem sok információt hordozó struktúra (járműindex + irány). A \texttt{StateMove} egy teljes átmenetet rögzít: forrásállapot, célállapot és az előbb említett lépést. Mivel a modellünkben minden lépésnek meg lehet tenni a fordítottját is, ezért egyenlőségvizsgálatnál nem teszünk különbséget a fordított lépéssel: $(A \to B) = (B \to A)$, így irányítatlan gráfot alkotunk a modellünkhöz.
\end{description}

\begin{figure}[H]
\centering
\includegraphics[width=0.8\textwidth]{model_solvers}
\caption{A keresőalgoritmusok osztálydiagramja}
\label{fig:model-solvers-class}
\end{figure}


\begin{description}
    \item[\texttt{RHSolver} hierarchia] Az absztrakt \texttt{RHSolver} alaposztály definiálja az eseményvezérelt interfészt és a közös logikát. Konkrét alosztályok:
    \begin{compactitem}
        \item \texttt{TabuSolver}: fenntart egy listát a legutóbbi állapotokról; kiválasztja a legjobb nem-tabu szomszédot.
        \item \texttt{BacktrackingSolver}: rendezett opciólisták vermét tartja fenn; visszalép zsákutcák esetén.
        \item \texttt{AcGraphSolver}: $f = g + h$ szerint rendezett \texttt{PriorityQueue}-t használ; egy szótárban követi a felfedezett állapotokat.
        \item \texttt{BFSDiscoverer} / \texttt{DFSDiscoverer}: a \texttt{Discoverer} (a gráfkeresést veszi alapul ezért a \texttt{AcGraphSolver}-ből öröklődik) osztály nem ellenőrzi és terminál megoldott állapot kiterjesztése esetén, a szélességi és mélységi felfedező alosztályai a kiértékelő függvényt definiálják felül.
    \end{compactitem}
    
    \item[\texttt{Heuristic} hierarchia] Absztrakt osztályhierarchia, amely a típusrendszeren keresztül kényszeríti ki a heurisztika tulajdonságait: \texttt{Heuristic} $\supset$ \texttt{NonNegativeHeuristic} $\supset$ \texttt{AdmissibleHeuristic} $\supset$ \texttt{MonotoneHeuristic}. Az A$^C$ gráfkereső fordítási időben megköveteli a \texttt{MonotoneHeuristic} típust.
\end{description}

Még érdemes megemlíteni a modellhez tartozó \texttt{Levels} Statikus osztályt, amely a szinteket leíró sztringeket és egy feldolgozó metódust tartalmaz, amely \texttt{RHGameState} példányokká alakítja azokat. Ezzel könnyen leírható egy pálya anélkül, hogy az objektumok pozíció struktúráját kéne megadni. A szintek hivatalos videójáték feladványaiból lettek válogatva.~\cite{rushhour_video_game}


\subsection{Nézet}
\label{sec:view}

A nézet kiemelt részének osztálydiagramja a \ref{fig:view-gameboard-class}. és \ref{fig:view-stategraph-class}. ábrákon látható.

\begin{figure}[H]
\centering
\includegraphics[width=1.0\textwidth]{view_gameboard}
\caption{Az játéktábla megjelenítéséhez használt osztályok osztálydiagramja}
\label{fig:view-gameboard-class}
\end{figure}

\begin{description}
    \item[\texttt{GameBoard}] Egy \texttt{RHGameState} megjelítéséért felelős objektum a szintbetöltés fülön a feladványok előnézetei ezt használják. Az játéktábla fülön egy \texttt{MainGameBoard} alosztály van példányosítva, ami ezeli a Kézi/Algoritmus mód váltását, a táblaállapot megjelenítését és a járművek nyilaiért felel. Feliratkozik a csúcskattintási eseményekre a manuális navigációhoz.
    \item[\texttt{VehicleNode} / \texttt{CarNode} / \texttt{BusNode}] 2D sprite-ok, amik a járműveket reprezentálják, egy GameBoard objektumhoz tartoznak (A Godot scene tree gyerek/szülő elérő metódusokon keresztül tudunk hivatkozni a kompozíció tagjaira). Ha a \texttt{MainGameBoard} hoz tartozik, akkor \texttt{Arrow} léptető nyíl objektumok  tartoznak hozzá a kézi mozgatáshoz. A nyilak csak egyszer vannak példányosítva, és ha nem lehet a nyíl által mezőre lépni, akkor az \texttt{IsActive} mező alapján ki lesz kapcsolva a láthatósága és interakciója.
\end{description}



\begin{figure}[H]
\centering
\includegraphics[width=1\textwidth]{view_stategraph}
\caption{Az állapotgráf megjelenítéséhez használt osztályok osztálydiagramja, a nyolcasfa megvalósítása a projekthez szabva}
\label{fig:view-stategraph-class}
\end{figure}




\begin{description}
    \item[\texttt{GraphScene}] A 3D gyökér-node, amely futtatja a szimulációt, meghívja a kirajzolást, és kezeli a nehéz számítással járó lassulást. Minden fizikai képkockában újjáépíti a Barnes–Hut nyolcasfát a csúcsok aktuális helyzete szerint. Minden fizikai képkockánál minden él és csúcs erőkiszámítását futtatja párhutamosan, majd a csúcsok elmozdulását is párhuzamosítva kiszámítja. A csúcsok kattintásához lefuttat egy ray castot a kamera aktuális állása szerint.
    \item[\texttt{Vertex}] Egy állapotot jelöl, egy körként van megjelnítve az állapotgráfon. Szálbiztos adatstruktúrát biztosít a rá gyakorolt erők akkumlálására. Prioritás-alapú effekt-rendszert tart fenn a színhez és mérethez. A felépített nyolcasfán kereszül approximálja a rá ható erőt.
    \item[\texttt{Edge}] Egy lépést jelöl, egy vonalként van megjelenítve az állapotgráfon. Az \texttt{ApplySpringForce} metódussal lehet a csúcsaira hozzáadni a rugóerőt. Lekérdezhető a színe a beállított effektusa alapján.
    % find a source that demonstrates godot MultiMeshInstance3D usage
    \item[\texttt{VertexDrawer, EdgeDrawer}] A renderelésért felelős osztályok, ahelyett, hogy minden csúcs él, külön lenne \texttt{Sprite3D} vagy \texttt{ImmediateMesh} objektum lenne, összegyűjtjük a komponenseket egy-egy tömbbe és azt felhasználva egyszerre rajzoljuk ki az összes csúcsot, élt, ezzel jelentősen felgyorsítva a renderelést.
    \item[\texttt{OctTreeNode}] A \textit{NetOctree} repository~\cite{octree_mcserep} alapján lett implementálva, csak a projekthez szükséges metódusokkal, és kiegészítve erőapproximálással. Implementálja a Barnes–Hut algoritmust három dimenziós térre.
\end{description}


Még említésre való osztályok a nézet rétegben:
\begin{description}
\item[\texttt{Camera3d}] A \texttt{GraphScene}-hez tartozik ortogonális projekciójú 3D kamera mozgással, forgatással és nagyítással az állapotgráf részletesebb megtekintéséhez. A \textit{3D-RTS-Camera} repository alapján készült. ~\cite{3d_rts_camera}
\item[\texttt{AlgoPlayer}] Vezérli az algoritmus futásának időzítését, feliratkozik/leiratkozik a solver eseményeiről, kezeli a lejátszást/szünetet/léptetést/újraindítást.
\item[\texttt{SolverSettingsTab}] Felhasználói felület az algoritmus- és heurisztika-választáshoz, tabu mérethez, véletlenszerűséghez és a BFS/DFS felfedezéshez.

\end{description}

\subsection{Adattárolás}

Az alkalmazás nem használ adatbázist. A rejtvények szintjei beégetett sztring literálokként tárolódnak a \texttt{Levels} osztályban. Minden szint egy 7 soros sztring: az első sor a szint címe, a maradék hat sor pedig a $6\times6$-os rácsot reprezentálja. Ez a reprezentáció kompakt, ember által olvasható és egyszerűen interpretálható.

\subsection{Keresőalgoritmusok és heurisztikák}
\label{sec:algorithms}

Ez a szakasz részletezi a három keresőalgoritmust és az irányításukhoz használt heurisztikus függvényeket.

\subsubsection{Heurisztikus függvények}

Mind a négy heurisztika megfelel három alapvető tulajdonságnak:
\begin{compactitem}
\item \textbf{Nemnegativitás}: $h(s) \geq 0$ minden $s$ állapotra.
\item \textbf{Megengedhetőség (Admissibility)}: $h(s) \leq h^*(s)$, ahol $h^*(s)$ a megoldáshoz szükséges minimális lépésszám $s$-ből.
\item \textbf{Monotonitás} (konzisztencia): bármely $s$ állapotra és annak egyetlen lépéssel elérhető $s'$ utódjára $| h(s) - h(s') | \leq 1$.
\end{compactitem}

Ezek a tulajdonságok szükségesek, hogy a gráfkeresés egyben A$^C$ gráfkeresés legyen, mivel ekkor egy állapot első kiterjesztésekor a hozzá vezető legrövidebb utat már megtaláltuk, többet nincs szükség kiterjeszteni azt az állapotot.


\paragraph{Null heurisztika.} $h(s) = 0$ minden állapotra. Nem ad információt az állapotról; a keresés nem informált brute-force felfedezéssé válik. Hasznos az alap algoritmus-viselkedés megfigyeléséhez.

Az alábbi heurisztikák mind speciálisan a Rush Hour $6 \times 6$-os táblájára és a 3. sor jobb szélén lévő kijáratára lettek implementálva. 

\paragraph{Távolság (Distance) heurisztika.} $h(s) = w_{\text{tábla}} - 1 - x_{\text{fő}}$, ahol $w_{\text{tábla}}$ a tábla szélessége $x_{\text{fő}}$ a főautó elejének oszlopa (0-tól indexelve). 6 mező széles táblára ez az érték 0 (megoldva) és 4 között mozog. Ez megengedhető, mivel a főautónak legalább ennyi mezőt kell elmozdulnia a kijáratig. Monoton, mert egyetlen lépés legfeljebb 1-gyel változtatja meg a főautó távolságát.

\paragraph{Távolság + torlaszok heurisztika.} A távolság heurisztikához hozzáad egy-egy költséget minden olyan jármű által elfoglalt mezőért, amely a főautó és a kijárat között van a kijárat sorában. Megengedhető, mert minden blokkoló elem eltávolítása legalább egy plusz lépést igényel. Monoton, mert egy lépés legfeljebb 1-gyel módosítja a távolságot vagy a torlaszok számát.

\paragraph{Legkevesebb lépés (Least-moves) heurisztika.} Kiterjeszti a távolság heurisztikát a blokkoló járművek eltávolítási költségének alsó becslésével. Minden, a főautó útját elálló járműre kiszámítja a minimális lépésszámot, amellyel félretolható, figyelembe véve: (a) az irányt, amerre tolni kell (buszoknál mindig lefelé a 4–6. sorokba, mivel a kijárat a 3. sorban van), és (b) a tolást akadályozó további járműveket. Ez az elemzés a monotonitás megőrzése érdekében már nem lehet ugyanígy kiterjeszteni a fő autót blokkoló autóknak a blokkolóinak eltolánsának lépésszámára, mivel ugyanezzel a becslési módszerrel már nem lenne monoton a heurisztika, egy ilyen autó elmozdítása 2-t is csökkenthet a heurisztikus függvény értékén. 

\subsubsection{Tabu keresés}
\label{sec:tabu}

% TODO cite glover1989tabu? ~\cite{glover1989tabu}
A tabu keresés egy lokális keresési meta-heurisztika. Minden lépésben felfedezi az összes szomszédos állapotot, kiszűri a tabu listán szereplőket, kiértékeli mindegyikre a heurisztikus függvényt és a heurisztika szerinti legjobb maradék szomszédra lép.

A tabu lista az utolsó $n$ meglátogatott állapotot tárolja. Főbb konfigurációk:
\begin{compactitem}
\item $n = 0$: mohó keresés memória nélkül, mindig a heurisztika szerinti legjobb szomszédra lép.
\item $n = 1$: hegymászó algoritmus; csak a közvetlenül megelőző állapotba való visszatérést kerüli el.
\item $n > 1$: elkerüli az utolsó $n$ állapotot. Megoldás nélkül leáll, ha minden szomszéd benne van a tabuban.
\end{compactitem}

A tabu keresés \emph{nem} garantálja a megoldás megtalálását, sem annak optimalitását. Ez egy gyors, lokális módszer, amely csapdába eshet az állapottér olyan régióiban, ahol kedvezőek a heurisztikus pontszámok, de nincs megoldás.

\subsubsection{Visszalépéses (Backtracking) keresés}
\label{sec:backtracking}

A visszalépéses keresés minden állapothoz a szomszédos még az úton még azelőtt nem látogatott állapotok heurisztika szerint rendezett sorozatát tartja fenn: minden lépésben a jelenlegi út utolsó állapotának a legjobb még nem próbált opciót vizsgálja meg. Ha már nincs új opció akkor leveszi az állapotot feléített út végéről, azaz algoritmus visszalép.


% TODO cite from https://people.inf.elte.hu/gt/mi/
\textbf{Leállási garancia}: a backtracking algoritmus véges állapotú állapotgráfra garantáltan leáll. Ha megoldás nélkül áll le, akkor a kezdőállapotból nem érhető el megoldás. Ugyanakkor csak az aktuális útvonalon lévő állapotokat kerüli el, a korábban már bejártakat nem. Ez azt jelenti, hogy többször is feltárhat nagy állapotcsoportokat, ami jelentős redundáns munkához vezethet.

% TODO cite cite this
\textbf{Időbonyolultság}: legrosszabb esetben a backtracking $O(V!)$ útvonalat jár be, ahol $V$ az elérhető állapotok száma. Pontosabban, minden egyszerű utat feltárhat a kezdőállapotból. Egy $b$ elágazási tényezőjű és $V$ állapotú térben a legrosszabb eset $O(b^V)$. A Rush Hour esetében egy adott feladványra az elágazási tényező legfeljebb $2 N_{\text{jármű}}$ ahol $N_{\text{jármű}}$ az autók és buszok száma összesen. Megfelelő heurisztikával az algoritmus a legtöbb esetben ennél sokkal gyorsabban talál megoldást, amennyiben elérhető.

\subsubsection{A\textsuperscript{C} gráfkeresés}
\label{sec:graphsearch}

% TODO cite A* ?
Az A$^C$ gráfkeresés egy informált keresőalgoritmus, amely az $f(n) = g(n) + h(n)$ kiértékelő függvényt használja, ahol $g(n)$ a mélység, $h(n)$ pedig a heurisztikus pontszám.

Mivel minden alkalmazható heurisztika nemnegatív, megengedhető és monoton, a következő tulajdonság teljesül: amikor egy állapotot kiveszünk a prioritási sorból kiterjeszteni, az az útvonal, amelyen elértük, a legrövidebb út az adott állapothoz. Ezért minden állapotot legfeljebb egyszer kell kiterjeszteni (az A$^C$-ben a „C” a „closed list”-re utal – a már kiterjesztett állapotokat nem terjesztjük ki újra).

Ez kiküszöböli a backtracking körbejárási és újrafelfedezési problémáit, de az algoritmus sok olyan állapotot is feltárhat, amelyek nem szükségesek a megoldáshoz, mivel megvizsgál minden állapotot, amelynek $f$-értéke kisebb az optimális megoldás költségénél.

\textbf{Időbonyolultság}: Az A$^C$ gráfkeresés minden állapotot legfeljebb egyszer terjeszt ki, ami legfeljebb $O(V)$ kiterjesztést jelent. Minden kiterjesztés legfeljebb $O(b)$ utódot generál ($b$ az elágazási tényező, max $2 \times N_{\text{jármű}}$), és minden beillesztés a prioritási sorba $O(\log V)$ költségű. Összesen: $O(V \cdot b \cdot \log V)$. A tárigény $O(V)$ a felfedezett és kiterjesztett állapotok szótára miatt.

\subsubsection{Barnes–Hut algoritmus és a nyolcasfa}
\label{sec:barneshut}

% TODO ~\cite{fruchterman1991graph}
Az erőalapú gráfelrendezés megköveteli a taszítóerők kiszámítását minden csúcspár között, ami képkockánként $O(V^2)$. Több ezer csúcsnál ez teljesítménybeli szűk keresztmetszetté válik.

A Barnes–Hut algoritmus~\cite{barneshut1986} ezt $O(V \log V)$-re csökkenti egy térbeli hierarchikus adatszerkezet – 3D-ben egy \textbf{nyolcasfa (octree)} – használatával. Az algoritmus menete:

\begin{enumerate}
\item \textbf{Nyolcasfa felépítése}: kiszámítjuk az összes csúcsot befoglaló tengelypárhuzamos kockát, azaz amibe minden csúcs belefér, ezt vesszük alapul és sorra beillesztjük a csúcsokat. Minden belső csomópont tárolja a részfájában lévő csúcsok számát és tömegközéppontját. A levelek egyetlen csúcsot tartalmaznak. A felépítés $V$ csúcsra $O(V \log V)$ időt igényel.\cite{barneshut1986}
\item \textbf{Erők kiszámítása}: minden $v$ csúcsra meghívjuk a fa gyökerére az erőkiszámítást. Minden $o$ csomópontnál:
    \begin{compactitem}
        \item Ha $s / d < \theta$ (ahol $s$ az $o$ kockájának oldalhosszúsága, $d$ pedig $v$ távolsága $o$ tömegközéppontjától), akkor az egész cellát egyetlen testként kezeljük, amelynek tömege megegyezik a benne lévő csúcsok számával. Ez a \emph{Barnes–Hut kritérium}.
        \item Ellenkező esetben rekurzívan bejárjuk $o$ csomópont gyerekeit, minden gyerekre kiszámítva/approximálva a $v$ csúcsra ható erőt és ezeket applikáljuk a $v$ csúcsra
    \end{compactitem}
    Az eljárás egy csúcsra átlagosan $O(\log V)$ időt igényel.~\cite{barneshut1986} Minden csúcsra összesen az erőkiszámítás is $O(V \log V)$ időt igényel.
% TODO find a citation for theta
\item A tipikusan $0.5 - 1.5$ intervallumra beállított $\theta$ paraméter szabályozza a pontosság és a sebesség közötti kompromisszumot. Az alkalmazás a $\theta = 1.0$ értéket használja.
\end{enumerate}

A két pont közötti taszítóerő a fordított négyzetes törvényt követi: $\vec{F} = -k \cdot m / d^2 \cdot \hat{r}$.

A nyolcasfát minden fizikai frissítésnél elölről felépítjük, mivel a csúcsok pozíciója folyamatosan változik. A faépítés $O(V \log V)$, az erő számítása csúcsonként $O(\log V)$, így a teljes időkomplexitás $O(V \log V)$ frissítésként.

%% ============================================================
\section{Implementáció}
%% ============================================================

\subsection{Felhasznált technológiák}

\begin{compactitem}
\item \textbf{Godot Engine 4.6} (Mono/.NET variáns)~\cite{godot_engine}: nyílt forráskódú játékmotor a 2D/3D rendereléshez, fizikához és felhasználói felülethez.
\item \textbf{C\# és .NET SDK 10.0}~\cite{dotnet10}: az összes alkalmazáslogika ezen a nyelven készült. Az egyike a Godot által támogatott nyelvek közül és ideális az objektum orientált logikához.
\end{compactitem}

Kihasznált godot funkciók:
\begin{compactitem}
\item \texttt{MultiMeshInstance3D} A csúcsok és élek batchelt rendereléséhez.
\item \texttt{Container} objektumok alkalmazása a felhasználói felület felépítéséhez. Beépítetett gombra, \texttt{CollisionShape2D}-re kattintási események használata.
\item \texttt{\_PhysicsProcess} metódus futtatása háttérszálon, hogy az interakciót ne blokkolja.
\end{compactitem}

\subsection{Eltérések az eredeti tervtől}

\paragraph{Erő kiértékelés és alvási optimalizáció.}
A fejlesztés során először minden csúcs és él egy saját megjelenítéssel rendelkező godot objektum volt.
A több ezer él és csúcs külön meghívott megjelenítése minden képkockán rendkívül számításigényesnek bizonyult. Ennek megoldására először az alábbi megoldás került implementációra:

\begin{compactenum}
\item Ahelyett, hogy az erőket közvetlenül a Godot fizikai API-ján keresztül alkalmaznánk, az erők egy csúcsonkénti (\texttt{Vector3 \_pendingForces}) mezőben gyűlnek.
\item Minden fizikai frissítés végén egy \texttt{Processor} node végigmegy a csúcsokon meghívja a kiértékelő függvényüket.
\item Egy dinamikus küszöbérték dönti el, hogy a függőben lévő erő elég jelentős-e az alkalmazáshoz. Ez a küszöb a csúcs „életkorával” (\texttt{\_timeAwake}) és az összes csúcs számával skálázódik, így nagy, stabilizálódott gráfoknál a kis maradék erőket figyelmen kívül hagyja, a jelentős erőket pedig meghívja a Godot függvényeivel.
\item Ha az erők tartósan a küszöb alatt maradnak, a csúcsokra beállított sebesség tompítás miatt el fogják érni a nyugalmi állapotot, ekkor a motor „alvó” módba teszi a \texttt{RigidBody3D}-t, leállítva a mozgását és a \texttt{\_PhysicsProcess} hívását, ezután a csúcs újra felébredhet, ha egy élnek a rugóereje a küszöb fölé kerül és applikálva lesz az erő a csúcsra.
\item Azok az élek, amelyeknek mindkét végpontja alszik, kihagyják a vonal-mesh frissítését, jelentős renderelési terhet levéve a rendszerről.
\end{compactenum}

Ez a lépcsőzetes optimalizáció (erőküszöb $\to$ fizikai alvás $\to$ render-kihagyás) jelentős teljesítménynövekedést eredményezzet, amikor az állapotgráf már a nyugalmi helyzetébe rendeződött. Viszont a legtöbb többezer csúcsos állapotgráfra a nyugalmi helyzet elérése túl hosszú volt ahhoz, hogy aktuális lehessen ez az optimalizáció, ezért ezt teljesen újra kellett gondolni és átalakítani a csúcsokat és éleket nem godotból származó osztályokká, és a kirajzolásukat kiszervezni külön osztályokba. Ez az optimalizáció már aktívan mozgó állapotgráfra is sokkal gyorsabb renderelést eredményezett. Ezen felül még több ezer csúcs esetén, a fizikai frissítések számát levisszük $60/s \rightarrow 10/s$-ra, így nem kell az alkalmazásnak bevárnia minden UI frissítéshez az elvárt egy fizikai frissítés lefutását, hanem több UI képkocka lefuthat fizikai frissítések között, és így reszponzív tud maradni a felhasználói felület.

%% ============================================================
\section{Tesztelés}
%% ============================================================

\subsection{Egységtesztek}

A következő egységtesztek ellenőrzik a program helyességét:

\begin{center}
\begin{longtable}{ | p{0.05\textwidth} | p{0.30\textwidth} | p{0.55\textwidth} | }

\hline
\multicolumn{3}{|c|}{\textbf{Egységtesztek}}
\\ \hline
\textbf{\#} & \textbf{Teszteset} & \textbf{Ellenőrzés}
\\ \hline \hline
\endfirsthead

\hline
\textbf{\#} & \textbf{Teszteset} & \textbf{Ellenőrzés}
\\ \hline \hline
\endhead

\hline
\endfoot
\endlastfoot

U1
& Szint-parzolás
& Minden szint sikeresen betöltődik, a járművek száma és pozíciója egyezik az elvárttal.
\\ \hline

U2
& Táblarács konzisztencia
& Parzolás után a $6\times6$-os rács helyesen rendeli a cellákat a járművekhez vagy $-1$-hez.
\\ \hline

U3
& Lépésgenerálás
& A \texttt{GetPossibleMoves()} csak érvényes lépéseket ad vissza (nincs átfedés, nincs kilépés a tábláról).
\\ \hline

U4
& Állapot immutabilitás
& A \texttt{WithMove()} új állapotot ad vissza; az eredeti állapot változatlan marad.
\\ \hline

U5
& Állapot egyenlőség
& Két azonos járműpozíciójú állapot egyenlő; az eltérőek nem.
\\ \hline

U6
& Hash konzisztencia
& Az egyenlő állapotok hash kódja azonos; a különbözőeké (nagy valószínűséggel) eltérő.
\\ \hline

U7
& Megoldás detektálása
& Egy állapot akkor és csak akkor megoldott, ha a főautó pozíciója $(5, 2)$.
\\ \hline

U8
& Null heurisztika
& Minden állapotra 0-t ad vissza.
\\ \hline

U9
& Távolság heurisztika tartomány
& Az érték $[0, 4]$ között van; csak megoldott állapotra 0.
\\ \hline

U10
& Heurisztika megengedhetőség
& Minden heurisztikára: $h(s) \leq$ tényleges minimális lépésszám (ismert feladványokon mérve).
\\ \hline

U11
& Heurisztika monotonitás
& Minden $s$ állapotra és $s'$ utódra: $h(s) \leq h(s') + 1$.
\\ \hline

U12
& Tabu solver leállása
& Megfelelő tabu mérettel egy megoldható feladványon a solver vagy megoldást talál, vagy jelzi a leállást.
\\ \hline

U13
& Backtracking teljesség
& Megoldható feladványon a backtracking megtalálja a megoldást. Ha a feladványra megoldhatatlan jelzi azt.
\\ \hline

U14
& A$^C$ optimalitás
& Adott feladványra ismert optimális hosszúságú megoldással az A$^C$ monoton heurisztikával optimális hosszúságó megoldást talál.
\\ \hline

U15
& StateMove szimmetria
& $(A \to B)$ megegyezik $(B \to A)$-val az élek irányítatlansága miatt.
\\ \hline

U16
& BFS/DFS felfedezési szám
& Állapotlimit esetén a felfedezés pontosan a megadott számnál áll le.
\\ \hline

\caption{Egységtesztek a programhoz}
\label{tab:unit-tests}
\end{longtable}
\end{center}

Az egységtesztek futtatásához a projekt könyvtárában az alábbi parancsot kell futtatni:

\begin{lstlisting}
dotnet add package gdUnit4.api --version 5.1.0-rc1
\end{lstlisting}

Ezután a szerkesztőből futtathatóak az egységtesztek.

\begin{figure}[H]
    \centering
    \includegraphics[width=0.5\textwidth]{unit_tests}
    \caption{Egységtesztek sikeres futtatása a szerkesztőben}
    \label{fig:unit-tests}
\end{figure}


% TODO node tests?

\subsection{Végfelhasználói tesztesetek}

A következő tesztesetek a \ref{tab:userstory}. táblázat felhasználói történetein alapulnak.

\begin{center}
\begin{longtable}{ | p{0.05\textwidth} | p{0.20\textwidth} | p{0.20\textwidth} | p{0.20\textwidth} | p{0.20\textwidth} | }

\hline
\multicolumn{5}{|c|}{\textbf{Végfelhasználói tesztesetek}}
\\ \hline
\textbf{\#} & \textbf{Művelet} & \textbf{Előfeltétel} & \textbf{Bemenet} & \textbf{Elvárt eredmény}
\\ \hline \hline
\endfirsthead

\hline
\textbf{\#} & \textbf{Művelet} & \textbf{Előfeltétel} & \textbf{Bemenet} & \textbf{Elvárt eredmény}
\\ \hline \hline
\endhead

\hline
\endfoot
\endlastfoot

E1
& Szint betöltése
& Nem fut algoritmus
& „Level 1” gomb
& Tábla mutatja az 1. szintet; a gráfon egy kék csúcs van.
\\ \hline

E2
& Kézi lépés
& Szint betöltve
& Kattintás egy autó előre nyilára
& Az autó elmozdul; új csúcs és él jelenik meg a gráfon.
\\ \hline

E3
& BFS felfedezés
& Szint betöltve
& Max áll.: 50; BFS indítás
& A gráf feltöltődik max. 50 csúccsal; a megoldások zöldek.
\\ \hline

E4
& Felfedezés leállítása
& Felfedezés fut
& Megállítás (\textit{Pause}) gomb
& A folyamat megáll; a csúcsok megmaradnak.
\\ \hline

E5
& Tabu keresés futtatása
& Szint betöltve
& Tabu: 5, Távolság heur.; Play
& Algoritmus fut, narancs/sárga kiemelések; státusz: „Running”.
\\ \hline

E6
& Algoritmus léptetése
& Algoritmus nem terminált
& Léptetés (\textit{Step}) gomb
& Egy lépés lefut; a lépésszámláló nő.
\\ \hline

E7
& Megoldás megtalálása
& Futás megoldható szinten
& Kivárás
& Státusz: „Solved”; megoldási út zöld; megoldási hossz látható.
\\ \hline

E8
& Ugrás csúcsra
& Egy szint be van töltve, valamennyi csúcs fel van fedezve
& Bal klikk egy csúcsra
& Tábla Kézi módba vált az adott állapottal; a csúcs piros lesz.
\\ \hline

E9
& Kamera forgatása
& Gráf látható
& Jobb egérgomb + húzás
& A kamera forog a gráf körül.
\\ \hline

E10
& Nagyítás
& Gráf látható
& Egérgörgő
& A nézet közelít vagy távolít.
\\ \hline

E11
& Kamera mozgatása
& Gráf látható
& W/A/S/D billentyűk
& A kamera elmozdul az lenyomott gomb szerint.
\\ \hline

E12
& Csúcsok elrejtése
& Gráf látható
& Hide kapcsoló
& Minden fehér csúcs és él el lesz rejtve.
\\ \hline

E13
& Solver újraindítása
& Futott már az algoritmus
& Restart gomb
& Algoritmus alaphelyzetbe áll; algoritmus kiemelések és lépésszám törlődnek.
\\ \hline

E14
& Módváltás (Kézi/Algo)
& Szint betöltve
& Manuális/Algoritmus gomb
& A tábla a manuális vagy az algoritmus módba vált.
\\ \hline

\caption{Felhasználói történeteken alapuló tesztesetek}
\label{tab:user-tests}
\end{longtable}
\end{center}